% !TeX root = ../组合数学笔记.tex
% 由主文档通过 \input 载入；请编译主文档。
\setcounter{chapter}{8} % 保留原书章号。
\chapter{Catalan 数与整数分拆}
\label{ch:classical-structures}

本章研究两类重要对象。合法括号、受限路径和有序树共享递归分解，因此出现同一数列；整数分拆则忘掉部件顺序，适合用重数和图形表示。仅有前几项计数相同不足以证明两类对象等数，必须建立双射或共同的完整递推。

\section{合法括号、Dyck 路径与平面二叉树}

由 $n$ 个左括号和 $n$ 个右括号组成的字符串，若每个前缀中左括号不少于右括号，称为合法括号串。将左括号变成上升步 $(1,1)$，右括号变成下降步 $(1,-1)$，得到从 $(0,0)$ 到 $(2n,0)$ 且不低于横轴的\emph{Dyck 路径（Dyck path）}；按步骤还原括号给出逆映射。

非空合法括号串可唯一写成 $(A)B$，其中外层右括号是首个使左右括号数重新相等的位置，$A,B$ 各为合法括号串。把空串对应为一个叶子，非空串对应为一个根，其左右子树分别编码 $A,B$，便得到每个内部节点恰有两个有序孩子的满平面二叉树。大小 $n$ 对应内部节点数，左右次序属于结构，不能任意交换。

\section{Catalan 递推、生成函数与闭式}

令 $C_n$ 为含 $n$ 对括号的合法串数，$C_0=1$。上述唯一分解立即给出
\begin{equation}
 C_{n+1}=\sum_{i=0}^nC_iC_{n-i},\qquad
 C(x)=1+xC(x)^2,\qquad C(x)=\sum_{n\ge0}C_nx^n.
 \label{eq:catalan}
\end{equation}
前几项为 $1,1,2,5,14,42$。方程中常数项必须为 $1$，所以形式解为 $C(x)=(1-\sqrt{1-4x})/(2x)$，平方根取常数项为 $1$ 的唯一形式级数；分子常数项为零，故除以 $x$ 合法。另一个代数根不是普通幂级数。

\begin{theorem}[Catalan 公式]
对 $n\ge0$，$C_n=\frac1{n+1}\binom{2n}{n}$。
\end{theorem}
\begin{proof}
$n=0$ 时成立。对 $n\ge1$，不受高度限制且终点高度为零的路径有 $\binom{2n}{n}$ 条。对于曾到达高度 $-1$ 的坏路径，反转起点至首次到达 $-1$ 的全部上下步。这段原来下降步比上升步多一，反转后全路径有 $n+1$ 个上升步和 $n-1$ 个下降步，终点高度为 $2$。反之，任意终点为 $2$ 的路径必到达 $1$，反转其首次到达 $1$ 前的步骤便得到坏路径；两操作互逆。因此坏路径数为 $\binom{2n}{n-1}$。相减得
\[
 C_n=\binom{2n}{n}-\binom{2n}{n-1}
 =\frac1{n+1}\binom{2n}{n}.
\]
\end{proof}

这称为\emph{反射法（reflection principle）}。它的关键是指定首次触及边界的位置，保证反射可逆。凸 $(n+2)$ 边形的三角剖分也有 $C_n$ 种（$n\ge1$）：固定一条边，其所属三角形唯一地把剩余部分分成两个较小多边形；将空的边区域计为一种，所得递推与式\eqref{eq:catalan}相同。

\section{整数分拆与 Ferrers 图}

\begin{definition}
整数 $n\ge0$ 的\emph{分拆（partition）}是非增正整数序列 $\lambda_1\ge\cdots\ge\lambda_\ell>0$，满足 $\sum_i\lambda_i=n$；$n=0$ 时只有空分拆。总数记为 $p(n)$，其中 $p(0)=1$。各 $\lambda_i$ 称为部件。
\end{definition}

用左对齐的方格表示每个部件，第 $i$ 行有 $\lambda_i$ 格，所得图称为\emph{Ferrers 图（Ferrers diagram）}。交换行列得到共轭分拆 $\lambda'_j=|\{i:\lambda_i\ge j\}|$。转置两次恢复原图，且最大部件与部件个数交换，故“最大部件至多 $k$”与“部件个数至多 $k$”的分拆等数。

例如 $5$ 的分拆为 $5$、$4+1$、$3+2$、$3+1+1$、$2+2+1$、$2+1+1+1$、$1+1+1+1+1$，故 $p(5)=7$。整数合成保留部件顺序，共有 $16$ 个；集合分拆保留元素标签，计数又不同。三者不是同一公式的不同记号。

\section{分拆生成函数与基本恒等式}

部件 $j$ 的重数可以任取非负整数，对生成函数贡献 $1+x^j+x^{2j}+\cdots$，不同部件重数独立，因此
\begin{equation}
 \sum_{n\ge0}p(n)x^n=\prod_{j\ge1}\frac1{1-x^j},\qquad
 \prod_{j\ge1}(1+x^j)=\prod_{j\ge1}\frac1{1-x^{2j-1}}.
 \label{eq:partition-products}
\end{equation}
无穷乘积在形式意义下良定义，因为 $x^n$ 的系数只涉及 $j\le n$。第二式由 $1+x^j=(1-x^{2j})/(1-x^j)$ 逐系数约去偶数因子得到，说明不同部件分拆与只含奇数部件的分拆等数。

该结论还有保持总和的双射。对奇数部件 $u$ 的重数写二进制展开 $m_u=\sum_j\varepsilon_j2^j$，将其替换成所有 $\varepsilon_j=1$ 时的部件 $2^ju$。所得部件互不相同，因为每个正整数的奇数部分与二的指数唯一。逆向把每个部件 $2^ju$ 拆成 $2^j$ 个 $u$，即可恢复原分拆。例如 $3+3+1+1+1$ 对应 $6+2+1$。

\section{习题}
\begin{enumerate}
  \item\label{ex:classical-structures-1} 写出全部五个三对括号的合法串，给出其中一个对应的树和 Dyck 路径编码。
  \item\label{ex:classical-structures-2} 求凸七边形的三角剖分数，并说明参数为何不是 $C_7$。
  \item\label{ex:classical-structures-3} 求分拆 $(4,2,1)$ 的共轭，写出部件至多为 $3$ 的分拆生成函数。
  \item\label{ex:classical-structures-4} 用不同部件与奇数部件之间的双射，把 $8+6+3+1$ 转化为只含奇数部件的分拆，并核对总和。
\end{enumerate}

